\section{推理优化}
% xudong
为提升 Baichuan-M3 在医疗应用中的可用性与高效部署，我们实施了两类推理优化策略。
首先，为加速文本生成，引入 Gated Eagle-3 推测解码方法，显著提升推理吞吐。
其次，采用先进量化技术在不损失准确性的前提下降低模型内存需求。
这些推理优化降低了部署门槛，支持先进医学大模型的更广泛应用。


\subsection{推测解码}

推测解码是一种用于自回归生成的推理加速技术。它使用轻量草案模型提出多个候选 token，再由目标模型并行验证。算法接受最长的已验证前缀并丢弃其余候选，使目标模型在每次验证步骤中提交多个 token，从而提升解码吞吐。

在 Eagle-3~\cite{li2025eagle} 框架中，草案模型额外条件化于目标模型的隐状态以利用更丰富的语义信息。然而，目标模型与草案模型之间显著的容量差距可能引发表示失配：高维且信息密集的隐状态可能压垮轻量草案模型，削弱其有效利用辅助信号的能力，导致候选接受率降低，从而限制可达加速比。

为缓解该问题，我们在 Eagle-3 草案模型中引入 Gated-Attention~\cite{qiu2025gated} 模块（称为 Gated Eagle-3；见图~\ref{fig:gated_eagle3}），提供动态可学习的信息注入调节机制。具体而言，注意力输出通过门控单元，并与由当前层输入动态生成的门控向量进行逐元素乘法调制。该设计实现了细粒度、逐维的信息流控制，突出关键特征并抑制冗余或噪声成分。

如附录~\ref{app:eagle_3} 所示，Gated Eagle-3 的平均接受长度提升 0.31，平均吞吐相较 Eagle-3 基线提高 12\%。这些结果表明，对目标模型信息进行选择性调节可更可靠地转化为实际解码加速，并为后续改进提供参考。

\begin{figure}
    \centering
    \includegraphics[width=1.0\linewidth]{image/gated-eagle3.pdf}
    \caption{Gated Eagle-3 草案模型示意图。}
    \label{fig:gated_eagle3}
\end{figure}



\subsection{量化}
为在资源受限场景下降低 Baichuan-M3 的 GPU 内存占用，我们采用 INT4 权重量化。然而，专家混合（MoE）模型的稀疏激活特性给标准量化校准带来挑战~\cite{DBLP:journals/corr/abs-2503-21135}：校准语料通常只激活部分专家，导致校准偏置。结果是，频繁激活的专家量化准确，而较少激活的专家量化误差更大，可能在推理时引发不稳定且难以预测的精度退化。

为解决该问题，我们提出自生成校准方案以促进专家覆盖的均匀性。具体而言，我们构建多领域提示集合，并使用 BF16 模型生成高质量响应进行校准。该方法具有两点优势：（i）多样提示可激活几乎所有专家，为每个专家提供充足样本，缓解校准偏置；（ii）自生成响应促使量化模型匹配 BF16 模型的输出分布，从而减少分布差异。
量化参数训练基于 AutoRound~\cite{DBLP:conf/emnlp/ChengZSCHLL24} 框架，量化格式遵循 GPTQ~\cite{DBLP:journals/corr/abs-2210-17323} 标准。



实证结果表明，INT4 量化后的 M3 模型在主流基准上相较 BF16 版本实现近乎无损的性能。该结果验证了我们提出的 MoE 专用量化校准策略的有效性，并为大规模稀疏模型部署提供实践指导。